% Compile with pdfLaTeX. All article text and references are in this file.
% This is a revised working manuscript, not a CSSS-issued template.
\documentclass[11pt,letterpaper]{article}
\usepackage[T1]{fontenc}
\usepackage[utf8]{inputenc}
\usepackage{lmodern}
\usepackage[margin=0.9in,headheight=14pt,headsep=18pt]{geometry}
\usepackage{microtype}
\usepackage{parskip}
\usepackage{xcolor}
\usepackage{titlesec}
\usepackage{fancyhdr}
\usepackage{enumitem}
\usepackage{needspace}
\usepackage{xurl}
\usepackage{hyperref}
\definecolor{ink}{HTML}{193B4A}
\hypersetup{colorlinks=true,urlcolor=ink,linkcolor=ink,
 pdftitle={Beyond the Karman Line: Keeping Orbital Safety Plans Current},
 pdfauthor={Adhvaith Ananth and Renish Joshi},pdfsubject={Policy analysis; revised working manuscript}}
\urlstyle{same}
\setcounter{secnumdepth}{0}
\setlength{\parskip}{6pt}
\setlength{\emergencystretch}{2em}
\linespread{1.04}
\widowpenalty=10000
\clubpenalty=10000
\raggedbottom
\titleformat{\section}{\sffamily\bfseries\fontsize{26}{29}\selectfont\color{ink}}{}{0pt}{}
\titleformat{\subsection}{\sffamily\fontsize{16}{20}\selectfont\color{ink}}{}{0pt}{}
\titleformat{\subsubsection}{\sffamily\bfseries\large\color{ink}}{}{0pt}{}
\titlespacing*{\section}{0pt}{0pt}{7pt}
\titlespacing*{\subsection}{0pt}{0pt}{12pt}
\titlespacing*{\subsubsection}{0pt}{13pt}{5pt}
\setlist[enumerate]{leftmargin=1.5em,itemsep=6pt,parsep=0pt}
\providecommand{\tightlist}{\setlength{\itemsep}{0pt}\setlength{\parskip}{0pt}}
\pagestyle{fancy}
\fancyhf{}
\fancyhead[L]{\scriptsize\sffamily ORBITAL SAFETY AND SPACE POLICY}
\fancyhead[R]{\scriptsize\sffamily Adhvaith Ananth and Renish Joshi}
\fancyfoot[C]{\small\thepage}
\renewcommand{\headrulewidth}{0.3pt}
\begin{document}
\hypertarget{beyond-the-kuxe1rmuxe1n-line}{%
\section{Beyond the Kármán Line}\label{beyond-the-kuxe1rmuxe1n-line}}

\hypertarget{keeping-orbital-safety-plans-current}{%
\subsection{Keeping Orbital Safety Plans
Current}\label{keeping-orbital-safety-plans-current}}

\textbf{Adhvaith Ananth and Renish Joshi}\\
Policy analysis \textbar{} Revised working manuscript \textbar{} 4
October 2026

A satellite can still take pictures or relay signals after a fault makes
it harder to avoid a collision or leave orbit. When that happens, the
people responsible for its safety need an updated account of what it can
do.

\hypertarget{after-the-launch}{%
\subsubsection{After the launch}\label{after-the-launch}}

The Kármán line is a familiar way to describe the edge of space. The
Fédération Aéronautique Internationale uses a height of 100 kilometres
for astronautical records.{[}1{]} The Outer Space Treaty does not define
where outer space begins.{[}2{]} For satellites already in orbit,
though, another question deserves attention: how should oversight
respond when a spacecraft's condition changes?

Imagine an Earth-observation satellite whose camera still works but
whose propulsion system has become unreliable. Its operator had planned
to use that system to avoid close approaches and lower the orbit at the
end of the mission. The satellite can still provide a service, but its
safety plan now rests on equipment that may not work as expected.

This is a hypothetical example. It helps explain why an approval issued
before launch cannot be the only information used to judge a mission
years later. Engineers may understand the fault, while another operator
or a regulator is still working from an older description of the
spacecraft.

We propose a shared record of the capabilities that matter for orbital
safety. Operators would update it when a significant change affects
their ability to maneuver, respond to commands or carry out their
disposal plan. The discussion focuses on civil missions and the United
States. It draws on published sources; it does not measure how often the
reporting problem occurs or claim that the proposal has already been
tested.

\hypertarget{start-with-the-rules-already-in-place}{%
\subsubsection{Start with the rules already in
place}\label{start-with-the-rules-already-in-place}}

Article VI of the Outer Space Treaty makes states internationally
responsible for national space activities. It also requires
authorization and continuing supervision of nongovernmental activities.
The treaty leaves the choice of domestic regulator and reporting
procedures unspecified.{[}2{]}

In the United States, the FAA licenses commercial launch and reentry
operations within its jurisdiction. In March 2026, it completed the
transition to Part 450 after a period in which the older rules and the
new framework operated together.{[}3{]} That work addresses important
parts of a mission, while other authorities deal with activities such as
satellite communications.

FCC satellite rules already require information about debris, collision
risk, maneuverability and disposal. Section 25.171 requires operators to
keep emergency and interference contact details current.{[}4{]} Any new
reporting proposal should reuse this information where possible.

The existing disposal rules also recognise that a spacecraft's condition
matters. Section 25.283(e) sets a five-year maximum after mission end
for covered spacecraft in, or passing through, low-Earth orbit below
2,000 kilometres that plan uncontrolled atmospheric reentry. It defines
mission end by loss of collision-avoidance capability, or completion of
the primary mission for spacecraft without that capability.{[}4{]}

Our proposal concerns how changes are shared, including restrictions
that appear before a satellite loses a capability completely. Several
agencies having different responsibilities is not, on its own, evidence
of a problem. What matters is whether information reported to one
recipient reaches others who need it for a safety decision.

\hypertarget{be-clear-about-what-a-satellite-can-do}{%
\subsubsection{Be clear about what a satellite can
do}\label{be-clear-about-what-a-satellite-can-do}}

Calling a satellite ``maneuverable'' can leave important questions
unanswered. NASA describes several ways of reducing collision risk. Some
spacecraft use propulsion. Others can change their orbit through
differential drag, by turning to change how much area they expose to the
thin atmosphere. Certain spacecraft without trajectory-control
capability can turn to reduce the area facing an approaching
object.{[}5{]} These methods have different limits.

In our example, the faulty propulsion system might still support some
maneuvers. Engineers might also be unsure whether the problem is
temporary. A useful record would allow them to describe the restriction
and the uncertainty, instead of forcing a choice between ``working'' and
``failed.''

Every update should say when the assessment was made and what supports
it. If the operator is still investigating, that should be clear. The
record should also change when the fault is resolved, so that other
users do not keep treating a recovered satellite as impaired.

\hypertarget{what-the-record-would-contain}{%
\subsubsection{What the record would
contain}\label{what-the-record-would-contain}}

The record could be short. It would identify the spacecraft and its
operational contact, describe its current maneuvering limits, note any
restrictions on receiving or carrying out commands, and state whether
its disposal plan remains feasible. Supporting reports could be linked
where the recipients have permission to read them.

Access would need care. Another operator may need to know that a
satellite cannot maneuver during a particular period, without needing
its propulsion drawings or flight software. More detailed explanations
could be restricted to the operator and the authorities entitled to
receive them. Sharing a warning should not require publishing the whole
design.

There are existing formats to build on. In January 2026, the Office of
Space Commerce published updated TraCSS specifications for satellite
catalog information and orbital and conjunction messages.{[}6{]} TraCSS
is intended to provide basic space situational awareness services to
civil and private operators.{[}7{]} These formats show how information
can be exchanged consistently. They do not establish the shared
capability record proposed in this article.

A common record would also leave legal responsibilities in place.
Operators would remain responsible for their operations, and agencies
would act within their own powers. Making any new report compulsory
would require a legal basis and the applicable procedures. Connecting
systems would not give a traffic-information service the power to
approve or stop missions.

\hypertarget{decide-which-changes-need-reporting}{%
\subsubsection{Decide which changes need
reporting}\label{decide-which-changes-need-reporting}}

Reporting every minor fault would make the system difficult to use. We
suggest focusing on changes that affect collision avoidance, command
response or disposal. A camera outage with no effect on those functions
would generally fall outside the proposed record. A lasting propulsion
problem on a satellite that needs propulsion to leave orbit would fall
inside it.

The reporting deadline should reflect the situation. A brief loss of
contact may be routine for one mission and a serious concern for
another. A trial with operators could help establish useful categories
and deadlines. We do not propose a fixed number of hours without
evidence that it would work across those cases.

Operators should be able to submit an initial report before finishing
the investigation. They could state what they know, explain what remains
uncertain, and update it as testing continues. This would give
recipients useful information while allowing the engineering assessment
to develop.

A report would not automatically decide the satellite's future. In our
example, engineers might recommend recovery attempts, limits on
operations or earlier disposal while enough capability remains. Those
decisions need technical judgment. The purpose of the record is to make
sure the people involved are not relying on assumptions that have become
outdated.

\hypertarget{a-lesson-from-aeolus}{%
\subsubsection{A lesson from Aeolus}\label{a-lesson-from-aeolus}}

ESA's Aeolus wind-observation satellite reentered on 28 July 2023.
Before reentry, engineers used its remaining fuel for a series of
maneuvers that lowered its orbit. ESA called the operation an assisted
reentry and explained that Aeolus had not originally been designed to
perform those maneuvers.{[}8{]}

The example shows how much the end of a mission depends on the
spacecraft's remaining capabilities and the team's planning. It does not
prove that our reporting proposal would have changed the outcome, or
that other satellites could follow the same approach.

It does suggest a useful question for mission extensions: will keeping a
satellite in service make its eventual disposal less reliable? Operators
should be able to explain how an extension affects that plan. The fuel
or backup equipment worth preserving would depend on the mission; Aeolus
does not establish a single reserve suitable for every spacecraft.

\hypertarget{try-it-before-adding-a-requirement}{%
\subsubsection{Try it before adding a
requirement}\label{try-it-before-adding-a-requirement}}

A first trial could be voluntary and leave existing obligations
unchanged. Operators and agencies would map the information they already
exchange, then test whether one update could serve several authorized
uses. Small university missions should take part alongside larger
fleets. A system that only works for teams with dedicated compliance
staff would have limited value.

The trial should measure the work involved. How long does an update
take? Are operators receiving fewer duplicate requests? Do the right
people receive a warning, and can they understand it? A simple web form
could serve a small team, while an automated connection might suit a
large fleet.

The amount of evidence requested should depend on the possible
consequences and the uncertainty. A small spacecraft is not
automatically low risk. Equally, fleet size alone does not tell us how
well an operator manages safety. The review process should allow
operators to explain their circumstances and challenge a classification
they consider inaccurate.

\hypertarget{where-the-proposal-could-fall-short}{%
\subsubsection{Where the proposal could fall
short}\label{where-the-proposal-could-fall-short}}

Existing coordination may already meet much of this need. The trial
should compare the shared record with simpler improvements, such as
keeping contacts current or making existing reports easier to reuse. If
those changes work as well with less effort, they should be preferred.
This article does not establish that a new central database is
necessary.

The record would also depend on the quality of information submitted.
Named contacts, dates and supporting evidence would make reports easier
to check, but they would not guarantee accuracy. Procedures should
distinguish an honest early assessment from a misleading claim.
Reporting a fault should not, by itself, be treated as misconduct.

Security rules would need to specify who can read or change a record,
how long it is kept and how errors are corrected. Some warnings might be
shared widely while the technical details remain restricted. Whether
that arrangement provides enough information should be tested with the
teams expected to use it.

There are limits beyond the reporting system. A U.S. arrangement cannot
ensure participation by every foreign operator. International use would
need cooperation and compatible practices. A record also cannot repair
hardware, remove debris or settle liability after a collision. It could
improve the information available for decisions, but the decisions and
the physical work would still have to follow.

\hypertarget{keeping-the-plan-useful}{%
\subsubsection{Keeping the plan useful}\label{keeping-the-plan-useful}}

The Kármán line describes a boundary. Our concern is what happens during
the years a spacecraft spends beyond it. Components age, faults occur
and plans change. The information used for oversight should keep pace
with those changes.

A shared capability record is worth testing if it helps operators and
regulators act on the same, current understanding of a spacecraft. The
test should be practical: better information, fewer duplicate requests
and a reasonable burden on the teams doing the work. If the record adds
paperwork without improving decisions, it should be changed or dropped.

The goal is to keep a mission's safety plan useful after launch,
including when the spacecraft can no longer do everything its designers
expected.

\hypertarget{references}{%
\needspace{8\baselineskip}
\subsubsection{References}\label{references}}

\begingroup\small
\begin{enumerate}
\def\labelenumi{\arabic{enumi}.}
\item
  Sanz Fernández de Córdoba, S. \emph{100km Altitude Boundary for
  Astronautics}. Fédération Aéronautique Internationale, updated 21 June
  2004. \href{https://www.fai.org/page/icare-boundary}{FAI source}.
\item
  Kopal, V. \emph{Treaty on Principles Governing the Activities of
  States in the Exploration and Use of Outer Space, including the Moon
  and Other Celestial Bodies: Introductory Note}. United Nations
  Audiovisual Library of International Law, September 2008. Discusses
  Articles VI and VIII and the treaty's undefined terms.
  \href{https://legal.un.org/avl/ha/tos/tos.html}{UN source}.
\item
  Federal Aviation Administration. \emph{FAA Streamlines Commercial
  Space License Approvals}. 17 March 2026.
  \href{https://www.faa.gov/newsroom/faa-streamlines-commercial-space-license-approvals}{FAA
  source}.
\item
  Federal Communications Commission. \emph{47 CFR Part 25: Satellite
  Communications}, especially §§ 25.114(d)(14), 25.171 and 25.283(e).
  Electronic Code of Federal Regulations; consulted 4 October 2026.
  \href{https://www.ecfr.gov/current/title-47/chapter-I/subchapter-B/part-25}{Regulatory
  text}.
\item
  NASA. \emph{Step 3: Close Approach Risk Mitigation}. Conjunction
  Assessment Risk Analysis programme; consulted 4 October 2026.
  \href{https://www.nasa.gov/cara/step-3-close-approach-risk-mitigation/}{NASA
  source}.
\item
  Office of Space Commerce. \emph{OSC Publishes Updated TraCSS
  Specifications}. 22 January 2026.
  \href{https://space.commerce.gov/osc-publishes-updated-tracss-specifications/}{Commerce
  source}.
\item
  Office of Space Commerce. \emph{OSC Opens TraCSS ``Waitlist'' for
  Satellite Operators}. 4 February 2026. Cited for the service's stated
  purpose, not its current enrollment status.
  \href{https://space.commerce.gov/osc-opens-tracss-waitlist-for-satellite-operators/}{Commerce
  source}.
\item
  European Space Agency. \emph{Aeolus: a historic end to a trailblazing
  mission}. 29 July 2023.
  \href{https://www.esa.int/Newsroom/Press_Releases/Aeolus_a_historic_end_to_a_trailblazing_mission}{ESA
  source}.
\end{enumerate}
\endgroup

\hypertarget{about-the-authors}{%
\subsubsection{About the authors}\label{about-the-authors}}

Adhvaith Ananth and Renish Joshi are students at the Texas Academy of Mathematics and
Science at the University of North Texas. They both conduct undergraduate research at the University of North Texas. Their interests include
aerospace engineering, computational simulations, and public policy,
particularly how spacecraft design and operational decisions affect the
governance of space activities.

\end{document}